% GENERATED FILE. Edit canonical lessons, then run npm run generate:books.
% canonical-source-hash: fnv1a64:510bb95c545976c8
% canonical-lessons: MR-C74-duhkhi, MR-C74-vadil, MR-C74-ai, MR-C74-mulga, MR-C74-mulgi

\chapter{Sad, Father, Mother, Son, Daughter}
\label{ch:mr-sad-father-mother-son-daughter}
\hlchaptermodality{74}

\begin{chapteropening}
\textbf{By the end of this chapter:} \emph{I can say sad, a father, a mother, a son and a daughter.}

Complete the last lesson of chapter 74: i can say sad, a father, a mother, a son and a daughter.
\end{chapteropening}

\section[duḥkhī]{\mr{दुःखी} (duḥkhī) — sad}

\label{lesson:MR-C74-duhkhi}



\begin{quote}
Before the new one: say the Marathi for hunger, then the Marathi for happy.
\end{quote}

\subsection*{You'll want to know: \mr{दुःखी}}
\textbf{\mr{दुःखी}} — \emph{duḥkhī} — \textquotedblleft{}sad\textquotedblright{}.

\begin{grammarlens}[title={What you've built}]
One more word for this chapter.
\end{grammarlens}

\subsection*{Your turn}
\begin{itemize}
\raggedright
  \item \emph{Say it:} \emph{duḥkhī}
  \item \emph{Say it:} \emph{duḥkhī}, once more
  \item \emph{Say it:} say \emph{duḥkhī}
  \item \emph{Recall:} say the Marathi for hunger, then the Marathi for happy, then say \emph{duḥkhī} again
\end{itemize}

\begin{tcolorbox}[breakable,colback=teal!4,colframe=teal!35!black,title={Before you move on}]
What does \mr{दुःखी} mean? (\textquotedblleft{}Sad\textquotedblright{}.) Say it once more.
\end{tcolorbox}
\section[vaḍīl]{\mr{वडील} (vaḍīl) — a father}

\label{lesson:MR-C74-vadil}



\begin{quote}
Before the new one: say the Marathi for happy, then the Marathi for sad.
\end{quote}

\subsection*{You'll want to know: \mr{वडील}}
\textbf{\mr{वडील}} — \emph{vaḍīl} — \textquotedblleft{}a father\textquotedblright{}.

\begin{grammarlens}[title={What you've built}]
One more word for this chapter.
\end{grammarlens}

\subsection*{Your turn}
\begin{itemize}
\raggedright
  \item \emph{Say it:} \emph{vaḍīl}
  \item \emph{Say it:} \emph{vaḍīl}, once more
  \item \emph{Say it:} say \emph{vaḍīl}
  \item \emph{Recall:} say the Marathi for happy, then the Marathi for sad, then say \emph{vaḍīl} again
\end{itemize}

\begin{tcolorbox}[breakable,colback=teal!4,colframe=teal!35!black,title={Before you move on}]
What does \mr{वडील} mean? (\textquotedblleft{}A father\textquotedblright{}.) Say it once more.
\end{tcolorbox}
\section[āī]{\mr{आई} (āī) — a mother}

\label{lesson:MR-C74-ai}



\begin{quote}
Before the new one: say the Marathi for sad, then the Marathi for a father.
\end{quote}

\subsection*{You'll want to know: \mr{आई}}
\textbf{\mr{आई}} — \emph{āī} — \textquotedblleft{}a mother\textquotedblright{}.

\begin{grammarlens}[title={What you've built}]
One more word for this chapter.
\end{grammarlens}

\subsection*{Your turn}
\begin{itemize}
\raggedright
  \item \emph{Say it:} \emph{āī}
  \item \emph{Say it:} \emph{āī}, once more
  \item \emph{Say it:} say \emph{āī}
  \item \emph{Recall:} say the Marathi for sad, then the Marathi for a father, then say \emph{āī} again
\end{itemize}

\begin{tcolorbox}[breakable,colback=teal!4,colframe=teal!35!black,title={Before you move on}]
What does \mr{आई} mean? (\textquotedblleft{}A mother\textquotedblright{}.) Say it once more.
\end{tcolorbox}
\section[mulgā]{\mr{मुलगा} (mulgā) — a son; a boy}

\label{lesson:MR-C74-mulga}



\begin{quote}
Before the new one: say the Marathi for a father, then the Marathi for a mother.
\end{quote}

\subsection*{You'll want to know: \mr{मुलगा}}
\textbf{\mr{मुलगा}} — \emph{mulgā} — \textquotedblleft{}a son; a boy\textquotedblright{}.

\begin{grammarlens}[title={What you've built}]
One more word for this chapter.
\end{grammarlens}

\subsection*{Your turn}
\begin{itemize}
\raggedright
  \item \emph{Say it:} \emph{mulgā}
  \item \emph{Say it:} \emph{mulgā}, once more
  \item \emph{Say it:} say \emph{mulgā}
  \item \emph{Recall:} say the Marathi for a father, then the Marathi for a mother, then say \emph{mulgā} again
\end{itemize}

\begin{tcolorbox}[breakable,colback=teal!4,colframe=teal!35!black,title={Before you move on}]
What does \mr{मुलगा} mean? (\textquotedblleft{}A son; a boy\textquotedblright{}.) Say it once more.
\end{tcolorbox}
\section[mulgī]{\mr{मुलगी} (mulgī) — a daughter; a girl}

\label{lesson:MR-C74-mulgi}



\begin{quote}
Before the new one: say the Marathi for a mother, then the Marathi for a son; a boy.
\end{quote}

\subsection*{You'll want to know: \mr{मुलगी}}
\textbf{\mr{मुलगी}} — \emph{mulgī} — \textquotedblleft{}a daughter; a girl\textquotedblright{}.

\begin{grammarlens}[title={What you've built}]
One more word for this chapter.
\end{grammarlens}

\subsection*{Your turn}
\begin{itemize}
\raggedright
  \item \emph{Say it:} \emph{mulgī}
  \item \emph{Say it:} \emph{mulgī}, once more
  \item \emph{Say it:} say \emph{mulgī}
  \item \emph{Recall:} say the Marathi for a mother, then the Marathi for a son; a boy, then say \emph{mulgī} again
\end{itemize}

\begin{tcolorbox}[breakable,colback=teal!4,colframe=teal!35!black,title={Before you move on}]
What does \mr{मुलगी} mean? (\textquotedblleft{}A daughter; a girl\textquotedblright{}.) Say it once more.
\end{tcolorbox}
